\clearpage\pagestyle{otMain}\setlength{\otBlocEspacement}{2.5mm}\fontsize{11}{15}\selectfont\corrigesChapitre{13}{Le théorème de Pythagore}
\begin{blocLivre}{cor-13-1}\begin{corrige}{1}{Chercher l’hypoténuse}\demarche{Additionner les carrés des côtés de l’angle droit.}
Dans ABC rectangle en A, Pythagore donne $BC^2=6^2+8^2=100$. Comme BC est une longueur, $BC=\sqrt{100}=\res{10\ \mathrm{cm}}$. Elle dépasse bien chacun des deux autres côtés.
\label{fig-cor-13-1}\figureGeo[corrige]{pyth-01}
\end{corrige}\end{blocLivre}
\begin{blocLivre}{cor-13-2}\begin{corrige}{2}{Chercher un petit côté}\demarche{Soustraire au carré de l’hypoténuse.}
BC est l’hypoténuse : $13^2=5^2+AC^2$. Donc $AC^2=169-25=144$ et $\res{AC=12\ \mathrm{cm}}$. Contrôle : $25+144=169$.
\label{fig-cor-13-2}\figureGeo[corrige]{pyth-02}
\end{corrige}\end{blocLivre}
\begin{blocLivre}{cor-13-3}\begin{corrige}{3}{Une racine non entière}\demarche{Distinguer valeur exacte et approximation.}
$BC^2=49+81=130$, donc $\res{BC=\sqrt{130}\ \mathrm{cm}}$. Au dixième, $BC\approx\res{11{,}4\ \mathrm{cm}}$. On conserve la racine pour la valeur exacte.
\label{fig-cor-13-3}\figureGeo[corrige]{pyth-03}
\end{corrige}\end{blocLivre}
\begin{blocLivre}{cor-13-4}\begin{corrige}{4}{L’hypoténuse donnée}\demarche{Repérer l’hypoténuse avant le calcul.}
$AC^2=BC^2-AB^2=289-225=64$. Donc $\res{AC=8\ \mathrm{cm}}$. BC reste le plus long côté.
\label{fig-cor-13-4}\figureGeo[corrige]{pyth-04}
\end{corrige}\end{blocLivre}
\begin{blocLivre}{cor-13-5}\begin{corrige}{5}{Une réciproque}\demarche{Comparer deux calculs distincts.}
BC est le plus long côté. $BC^2=225$ et $AB^2+AC^2=81+144=225$. L’égalité permet d’appliquer la réciproque de Pythagore : $\res{\text{ABC est rectangle en A}}$.
\label{fig-cor-13-5}\figureGeo[corrige]{pyth-05}
\end{corrige}\end{blocLivre}
\begin{blocLivre}{cor-13-6}\begin{corrige}{6}{Un triangle non rectangle}\demarche{Utiliser l’inégalité pour exclure le caractère rectangle.}
BC est le plus long côté. $7^2=49$, tandis que $4^2+5^2=41$. Les résultats sont différents : ABC n’est pas rectangle. Si le triangle l’était, le théorème imposerait leur égalité.
\label{fig-cor-13-6}\figureGeo[corrige]{pyth-06}
\end{corrige}\end{blocLivre}
\begin{blocLivre}{cor-13-7}\begin{corrige}{7}{Diagonale d’un rectangle}\demarche{Faire apparaître un triangle rectangle.}
La diagonale est l’hypoténuse d’un triangle rectangle dont les deux autres côtés mesurent 8 et 6 cm. Son carré vaut $64+36=100$ : elle mesure $\res{10\ \mathrm{cm}}$.
\label{fig-cor-13-7}\figureGeo[corrige]{pyth-07}
\end{corrige}\end{blocLivre}
\begin{blocLivre}{cor-13-8}\begin{corrige}{8}{Une échelle dans un modèle}\demarche{Modéliser la situation avant de calculer.}
Le mur et le sol forment un triangle rectangle ; l’échelle est l’hypoténuse. $h^2=2{,}5^2-0{,}8^2=5{,}61$, donc $h=\sqrt{5{,}61}\approx\res{2{,}37\ \mathrm m}$. Les deux mesures étaient bien en mètres.
\label{fig-cor-13-8}\figureGeo[corrige]{pyth-08}
\end{corrige}\end{blocLivre}
\begin{blocLivre}{cor-13-9}\begin{corrige}{9}{Une distance dans un repère}\demarche{Lire les écarts sur les axes puis utiliser Pythagore.}
AB et AC sont portés par les deux axes perpendiculaires, donc ABC est rectangle en A. $AB=6$ et $AC=8$ cm ; Pythagore donne $\res{BC=10\ \mathrm{cm}}$. On justifie l’angle droit par le repère orthonormé.
\label{fig-cor-13-9}\figureGeo[corrige]{pyth-09}
\end{corrige}\end{blocLivre}
\begin{blocLivre}{cor-13-10}\begin{corrige}{10}{Trouver le bon sommet}\demarche{Placer l’angle droit en face du plus long côté.}
Le plus long côté est AB. $AB^2=110{,}25$ et $AC^2+BC^2=39{,}69+70{,}56=110{,}25$. La réciproque montre que le triangle est rectangle au sommet opposé à AB, donc $\res{\text{en C}}$.
\label{fig-cor-13-10}\figureGeo[corrige]{pyth-10}
\end{corrige}\end{blocLivre}
\begin{blocLivre}{cor-13-11}\begin{corrige}{11}{Un losange}\demarche{Utiliser les demi-diagonales.}
Un quart du losange est un triangle rectangle de côtés de l’angle droit 5 et 12 cm. Le côté du losange vaut $\sqrt{25+144}=\res{13}$ cm. Son périmètre est $4\times13=\res{52}$ cm.
\label{fig-cor-13-11}\figureGeo[corrige]{pyth-11}
\end{corrige}\end{blocLivre}
\begin{blocLivre}{cor-13-12}\begin{corrige}{12}{Retrouver une aire}\demarche{Retrouver la dimension avant de calculer l’aire.}
Pour la longueur $L$, $L^2=13^2-5^2=144$, donc $L=12$ cm. L’aire vaut $12\times5=\res{60\ \mathrm{cm}^2}$. Ne pas confondre la diagonale avec une dimension du rectangle.
\label{fig-cor-13-12}\figureGeo[corrige]{pyth-12}
\end{corrige}\end{blocLivre}
\begin{blocLivre}{cor-13-13}\begin{corrige}{13}{Deux diagonales}\demarche{Traiter chaque triangle rectangle séparément.}
La première diagonale vaut $\sqrt{8^2+6^2}=10$ m. La seconde vaut $\sqrt{10^2+24^2}=26$ m. On additionne les longueurs, pas leurs carrés : total $\res{36\ \mathrm m}$.
\label{fig-cor-13-13}\figureGeo[corrige]{pyth-13}
\end{corrige}\end{blocLivre}
\begin{blocLivre}{cor-13-14}\begin{corrige}{14}{Une précision décisive}\demarche{Comparer les carrés pour une décision exacte.}
La distance à couvrir vaut $\sqrt{4^2+7^2}=\sqrt{65}$ cm. Or $8^2=64<65$, donc $\sqrt{65}>8$. La barre est trop courte. Un arrondi trop grossier à 8 cm ferait perdre l’information décisive.
\label{fig-cor-13-14}\figureGeo[corrige]{pyth-14}
\end{corrige}\end{blocLivre}
\begin{blocLivre}{cor-13-15}\begin{corrige}{15}{Choisir un trajet}\demarche{Calculer les deux parcours avant de comparer.}
Le trajet sur les côtés vaut $12+5=17$ m. La diagonale vaut $\sqrt{12^2+5^2}=13$ m par Pythagore. Le gain est $17-13=\res4$ m. La comparaison porte sur des longueurs dans la même unité.
\label{fig-cor-13-15}\figureGeo[corrige]{pyth-15}
\end{corrige}\end{blocLivre}
